% Ad Familiares 9.21 (ad-familiares-09-21) --- English translation
% Translated by Claude (Anthropic) from the Latin (Perseus).
% Style guide: STYLE.md.

\ciceroLetterOpener{CICERO PAETO S.}

\ciceroSection{1} Do you really mean it? You think yourself out of your mind because you imitate, as you write, the ``thunderbolts'' of my words? You would be out of your mind if you could not do it; but as it is, since you actually outdo me, you should be laughing at me sooner than at yourself. So you have no need of that tag from Trabea --- rather, of a misfire \textit{[Greek: apoteugma]} of my own. But for all that, what do I look like to you in letters? Am I not talking with you in plebeian speech? After all, I do not always do it in the same style. For what has a letter in common with a court-speech or a public address? Why, even our court-cases we do not handle in any one way. Private suits, and those of slight stake, we conduct in a finer-spun style; cases of life or reputation more ornately, of course. Letters, on the other hand, we usually weave from everyday words.

\ciceroSection{2} But for all that, my dear Paetus, what came into your head to deny that any Papirius was ever anything but plebeian? For there were patricians of the lesser clans, of whom the chief was L. Papirius Mugillanus, who was censor with L. Sempronius Atratinus, having earlier been consul with the same colleague, in the three hundred and twelfth year after the founding of Rome; though in those days you used to be called Papisii. After him there were thirteen who held the curule chair before L. Papirius Crassus, who was the first to leave off being called Papisius. This man was made dictator, with L. Papirius Cursor as his master of the horse, in the four-hundred-and-fifteenth year after the founding of Rome, and four years afterwards he was consul with K. Duilius. After him followed Cursor himself, a very much honoured man; then L. Masso, holder of the aedileship; from him came many Massones. Of all of these --- patricians, every one --- I should like you to have the ancestor-masks. After them the Carbones and the Turdi come into the line.

\ciceroSection{3} These were plebeians; whom I move that you despise; for, apart from this Gaius Carbo whom Damasippus killed, there was no Carbo who was a citizen of any use to the commonwealth. We knew Gnaeus Carbo, and his brother the buffoon; what was more worthless than they? About this present friend of mine, Rubria's son, I say nothing. There were three brothers in that line --- Gaius, Gnaeus, and Marcus Carbo. Marcus, prosecuted by P. Flaccus, was condemned, a great thief, out of Sicily; Gaius, with L. Crassus prosecuting, is said to have taken cantharides; he was both a seditious tribune of the plebs and was thought to have laid violent hands on Publius Africanus. As for the one who was put to death at Lilybaeum by our Pompey, no one in my judgement was more worthless. His father, indeed, prosecuted by M. Antonius, is reputed to have been acquitted by means of shoemaker's blacking. Wherefore I move you go back to the patricians; you see how the plebeians have been a graceless lot.
